% !TEX root = ../metatime_monograph.tex
\chapter{Quark Primes}
\label{ch:quark_primes}

Each quark flavour is assigned a prime number based on its position in
the generation structure:

\begin{equation}
(u,d,s,c,b,t) = (3,5,7,11,13,17).
\label{eq:quark_primes}
\end{equation}

These assignments are not arbitrary; they follow from an ordering principle
that encodes the flavour geometry.

\section{Assignment Principle}

The quark primes are the primes corresponding to the generation index:

\[
\begin{array}{c|c|c}
\text{Generation} & \text{Quarks} & \text{Prime} \\ \hline
1 & u, d & p_2 = 3,\ p_3 = 5 \\
2 & s, c & p_4 = 7,\ p_5 = 11 \\
3 & b, t & p_6 = 13,\ p_7 = 17
\end{array}
\]

Here \(p_n\) denotes the \(n\)-th prime number: \(p_1 = 2\), \(p_2 = 3\),
\(p_3 = 5\), \(p_4 = 7\), \(p_5 = 11\), \(p_6 = 13\), \(p_7 = 17\).
The assignments skip \(p_1 = 2\) (which is reserved for the binary structure
of the information constant) and start from \(p_2 = 3\).

\section{Flavour Geometry}

The prime numbers encode geometric relationships between flavours:

\begin{description}
\item[Isospin splitting:] The difference \(d - u = 5 - 3 = 2\) equals the
annihilation constant \(L_4\), encoding the up-down mass splitting.
\item[Strangeness scale:] The strange quark prime \(s = 7\) equals the
structural constant \(L_3\), establishing the strangeness scale.
\item[Generation spacing:] The prime gaps increase with generation:
\(|3-5| = 2\), \(|7-11| = 4\), \(|13-17| = 4\), reflecting the growing
hierarchy between successive generations.
\end{description}

\section{Sum and Difference Combinations}

Several combinations of quark primes appear in the sector formulas:

\[
\begin{array}{lcl}
s+u = 10, & s-u = 4, \\
s^2 - u^2 - d^2 = 49 - 9 - 25 = 15, & u+d+s = 15, \\
c = 11, & b = 13, & t = 17.
\end{array}
\]

\section{Physical Interpretation}

The quark primes serve as flavour labels on the Poincar\'{e} disk.  They
encode the action contributions of each flavour to the mass formulas.  The
metaphor is that of a periodic table: the primes are the atomic numbers of
flavour space, from which the properties of composite states (hadrons) are
derived by arithmetic combination rules.
